\documentclass{ctexart}
\pagestyle{empty} % 禁用页眉页脚，包括页码
\setlength{\parindent}{2em}
\usepackage[
  paperwidth=17cm,
  paperheight=25cm,
  margin=6pt
]{geometry}
\usepackage{titlesec}

% 修改 \section 字号为 14pt，加粗
\titleformat{\section}
  {\normalfont\bfseries\large} % 样式：\Large 可改为 \huge, \small 等
  {\thesection}{0em}{}         % 编号和间距设置
  \titlespacing{\section}{0pt}{10pt}{0pt}

\usepackage{multicol}  
\usepackage{amsmath}   

\begin{document}


\section*{二连杆末端的雅克比矩阵}

考虑平面二连杆机械臂，其连杆长度分别为 $l_1$ 和 $l_2$，关节角为 $\theta_1$ 和 $\theta_2$。末端在世界坐标系中的位置为
\begin{align*}
    x & = l_1 \cos\theta_1 + l_2 \cos(\theta_1+\theta_2) \\
    y & = l_1 \sin\theta_1 + l_2 \sin(\theta_1+\theta_2)
\end{align*}

雅克比矩阵定义为末端位置对关节变量的偏导矩阵：

\begin{equation*}
    J(\theta)=
    \begin{bmatrix}
        \frac{\partial x}{\partial \theta_1} & \frac{\partial x}{\partial \theta_2} \\
        \frac{\partial y}{\partial \theta_1} & \frac{\partial y}{\partial \theta_2}
    \end{bmatrix}
    =
    \begin{bmatrix}
        -l_1\sin\theta_1 - l_2\sin(\theta_1+\theta_2) &
        -l_2\sin(\theta_1+\theta_2)                     \\
        l_1\cos\theta_1 + l_2\cos(\theta_1+\theta_2)  &
        l_2\cos(\theta_1+\theta_2)
    \end{bmatrix}
\end{equation*}


\end{document}